\Askhsh{37192}{2}
\begin{ylh}{2}
    \item 2.4. Ρίζες Πραγματικών Αριθμών
\end{ylh}
Στον πίνακα της τάξης σας είναι γραμμένες οι παρακάτω πληροφορίες (προσεγγίσεις)
$\sqrt{2} \simeq 1,41,\sqrt{3} \simeq 1,73,\sqrt{5} \simeq 2,24,\sqrt{7} \simeq 2,64$
\begin{alist}
\item Να επιλέξετε έναν τρόπο ώστε να αξιοποιήσετε τα παραπάνω δεδομένα (όποια θεωρείτε κατάλληλα) και να υπολογίσετε με προσέγγιση εκατοστού τους αριθμούς $\sqrt{20}, \sqrt{45}$ και $\sqrt{80}$
\monades{12}
\item Αν δεν υπήρχαν στον πίνακα οι προσεγγιστικές τιμές των ριζών, πως θα μπορούσατε να υπολογίσετε την τιμή της παράστασης $\dfrac{3 \cdot \sqrt{20} + \sqrt{80}}{\sqrt{45} - \sqrt{5}}$
\monades{13}
\end{alist}

